\chapter{\rm\bfseries Experiment 2: Tilt Effect on Grasp Force}
\label{ch:exp2}

\section{Objective}

Characterise how off-axis gripper tilt perturbs the contact force, separately
at boundary and safe depths, for the most-studied variant
(\texttt{finray\_26}). Test whether the finray's symmetry about its $xz$
plane is preserved or broken under physical operation.

\section{Tilt-Aware Pose Geometry: Verification}
\label{sec:tilt-verify}

Before reporting force results, the implementation of
equation~\eqref{eq:flange-target-tilted} was verified empirically.
Table~\ref{tab:tilt-geometry} shows the measured flange XY shift during
the stationary grasp sample versus the geometric prediction for
\texttt{finray\_26} at four representative (depth, tilt) cells.

\begin{table}[ht]
\centering
\caption[Tilt geometry verification]{Measured vs predicted flange-XY offset
during the stationary grasp sample at four cells. Predicted offset from
equation~\eqref{eq:flange-target-tilted} with $z_{\text{TCP}} = 94.911$\,mm.}
\label{tab:tilt-geometry}
\small
\begin{tabular}{lccccc}
\toprule
Cell & Tilt cmd & Predicted $\Delta \mathbf{p}_\mathcal{F}$ (mm) &
       Measured $\Delta \mathbf{p}_\mathcal{F}$ (mm) & Residual & Mean $|\mathbf{F}|$ \\
\midrule
$-1.0$\,mm   & $+2.5^{\circ}$ y & $(+4.14, +0.00)$ & $(+4.35, -0.03)$ & $0.21$\,mm & $-8.00$\,N \\
$-1.0$\,mm   & $-2.5^{\circ}$ y & $(-4.14, +0.00)$ & $(-4.02, -0.04)$ & $0.13$\,mm & $-3.39$\,N \\
$+10.0$\,mm  & $+2.5^{\circ}$ y & $(+4.14, +0.00)$ & $(+5.35, -0.28)$ & $1.25$\,mm & $-36.26$\,N \\
$+10.0$\,mm  & $-2.5^{\circ}$ y & $(-4.14, +0.00)$ & $(-2.93, -0.29)$ & $1.21$\,mm & $-35.95$\,N \\
\bottomrule
\end{tabular}
\end{table}

At the boundary depth, the residual is $\leq 0.21$\,mm---essentially the
lateral positioning noise of the arm under light load. At the safe depth,
both positive and negative tilts exhibit a $+1.2$\,mm common-mode offset
in the same $+x$ direction, regardless of tilt sign. We interpret this
common-mode offset as evidence of finger compliance under load: the
$\sim 36$\,N axial press deflects the soft finray blade laterally by
approximately 1.2\,mm in a structurally preferred direction. The
differential offset (between $+y$ and $-y$ tilts) remains correctly
predicted by the geometric model in both depth conditions.

This verification confirms that the control-frame implementation of
equation~\eqref{eq:flange-target-tilted} is correct, and any asymmetries
observed in the contact force reflect physical gripper behaviour rather
than control system errors.

\section{Force vs Tilt Magnitude at Fixed Depth}

Figure~\ref{fig:tilt-mag-sweep} shows the stationary force at $0^{\circ}$,
$+2.5^{\circ}$, and $+5^{\circ}$ tilt about the symmetric $x$ axis, at
both boundary ($-0.5$\,mm) and safe ($+10$\,mm) depths.

\begin{figure}[ht]
\centering
\includegraphics[width=0.95\linewidth]{finray_26_tilt_sweeps.png}
\caption[Tilt sweep panels]{Tilt-magnitude (left) and signed tilt (right)
sweeps for \texttt{finray\_26}. Left panel: x-axis tilt (symmetric axis)
at boundary (blue) and safe (green) depths shows a strong tilt-effect at
boundary that disappears at safe depth. Right panel: y-axis tilt
(asymmetric axis) shows a clearly asymmetric response at boundary depth
but a symmetric, flat response at safe depth.}
\label{fig:tilt-mag-sweep}
\end{figure}

At the boundary depth, force decreases monotonically from $-10.8$\,N at
$0^{\circ}$ to $-5.1$\,N at $+5^{\circ}$---a $53\%$ reduction. At the safe
depth, the entire sweep stays within 1\,N of the $0^{\circ}$ baseline.
The data are summarised in Table~\ref{tab:tilt-mag}.

\begin{table}[ht]
\centering
\caption[Tilt magnitude effect]{Stationary force vs x-axis tilt magnitude
at two depth conditions for \texttt{finray\_26}.}
\label{tab:tilt-mag}
\begin{tabular}{lcccc}
\toprule
Depth & $0^{\circ}$ & $+2.5^{\circ}$ x & $+5.0^{\circ}$ x & \% reduction at $5^{\circ}$ \\
\midrule
$-0.5$\,mm (boundary) & $-10.82$\,N & $-6.08$\,N & $-5.09$\,N & $53\%$ \\
$+10$\,mm (safe)      & $-36.51$\,N & $-35.99$\,N & $-36.93$\,N & $-1\%$ \\
\bottomrule
\end{tabular}
\end{table}

\section{Asymmetric Axis: Signed y-Tilt Sweep}

Figure~\ref{fig:y-tilt-asym} shows the asymmetry ratio
$|F_z(+2.5^{\circ}_y)| / |F_z(-2.5^{\circ}_y)|$ as a function of depth,
on a logarithmic axis.

\begin{figure}[ht]
\centering
\includegraphics[width=0.85\linewidth]{finray_26_y_tilt_asymmetry.png}
\caption[Asymmetry vs depth]{Asymmetry ratio between $+2.5^{\circ}$ and
$-2.5^{\circ}$ y-tilts as a function of press depth for
\texttt{finray\_26}. The finray exhibits strong directional ``handedness''
at boundary depths but loses it in the soft plateau.}
\label{fig:y-tilt-asym}
\end{figure}

At $-1.0$\,mm, the ratio is $2.36\times$---force at $+2.5^{\circ}$
($-8.00$\,N) is more than double the force at $-2.5^{\circ}$ ($-3.39$\,N).
At $-0.5$\,mm, the ratio is $1.38\times$. At $+10$\,mm, the ratio is
$1.009\times$, statistically indistinguishable from 1.

Table~\ref{tab:y-tilt-asym} reports the absolute values along with
$0^{\circ}$ baselines where available.

\begin{table}[ht]
\centering
\caption[Y-tilt asymmetry]{Stationary force at signed y-tilts for
\texttt{finray\_26}, with $0^{\circ}$ baselines from experiment 1.}
\label{tab:y-tilt-asym}
\begin{tabular}{lcccc}
\toprule
Depth & $F_z(-2.5^{\circ}_y)$ & $F_z(0^{\circ})$ & $F_z(+2.5^{\circ}_y)$ &
        $|F_+ / F_-|$ \\
\midrule
$-1.0$\,mm  & $-3.39$\,N  & $-4.09$\,N  & $-8.00$\,N   & $2.36\times$ \\
$-0.5$\,mm  & $-10.01$\,N & $-12.20$\,N & $-13.86$\,N  & $1.38\times$ \\
$+10$\,mm   & $-35.95$\,N & $-36.51$\,N & $-36.26$\,N  & $1.009\times$ \\
\bottomrule
\end{tabular}
\end{table}

\section{Combined Picture: Three-Tilt Force-Depth Curves}

Figure~\ref{fig:three-tilts} integrates the depth and tilt sweeps for
\texttt{finray\_26}.

\begin{figure}[ht]
\centering
\includegraphics[width=0.95\linewidth]{finray_26_Fz_vs_depth_three_tilts.png}
\caption[Three-tilt depth curves]{Stationary force vs depth for
\texttt{finray\_26} at $0^{\circ}$, $+2.5^{\circ}$ x, and $+5^{\circ}$ x
tilts. Background shading marks the regime boundaries identified in
experiment 1. The three tilt conditions diverge at boundary depths and
converge in the soft plateau.}
\label{fig:three-tilts}
\end{figure}

The three lines diverge in the boundary regime and re-converge in the soft
plateau. Notably, the convergence is complete by $+10$\,mm; the soft
plateau is tilt-robust both in magnitude and in sign.

\section{Rotation Tracking Verification}

A natural concern is whether the Franka arm faithfully realises the
commanded tilt orientation. The 30\,Hz session log records actual and
commanded RPY at every sample, enabling a direct measurement of orientation
tracking error
\begin{equation}
  \delta(t) = \| \mathrm{RPY}_{\text{actual}}(t) -
                 \mathrm{RPY}_{\text{cmd}}(t) \|_2.
\end{equation}
Across five representative sessions, the mean steady-state $\delta$ was
$0.05^{\circ}$ with a worst-case $0.13^{\circ}$. Transient excursions up to
$\sim 5^{\circ}$ occur exclusively within the \texttt{orientation\_lock}
stage at session start, before any descent or grasp, and are filtered out
of trial-level analyses.

A $2.5^{\circ}$ commanded tilt is therefore $50\times$ the noise floor of
the arm's rotation tracking; the observed force differences cannot be
attributed to arm rotation error.
